\documentclass[11pt, a4paper]{article}
\usepackage[margin=2cm]{geometry}
\usepackage{enumitem}
\usepackage{booktabs}
\usepackage{hyperref}
\usepackage{graphicx}
\usepackage{tabularx}
\usepackage{longtable}
\usepackage{array}
\usepackage{titling}
\setlength{\parskip}{2pt}
\setlist{nosep, leftmargin=*}

\hypersetup{
    colorlinks=true,
    linkcolor=blue,
    urlcolor=blue
}

\title{\textbf{Assignment 3: Sprint, Git Workflow, MVP v1}\\[0.5em]
\large Software Project --- Spring 2026}
\author{}
\date{}

\begin{document}
\maketitle
\thispagestyle{empty}

\begin{center}
\begin{tabular}{ll}
\toprule
\textbf{Project} & PoetryForYou --- Telegram Poetry Learning Bot \\
\textbf{Customer} & Nursultan Askarbekuly \\
\textbf{Team Member} & Egor Zhukov (Solo Developer) \\
\bottomrule
\end{tabular}
\end{center}

\subsection*{What I Did This Sprint}
\begin{enumerate}
    \item Groomed the product backlog: re-estimated, added acceptance criteria, broke large items into sub-tasks.
    \item Created a GitHub Milestone for Sprint~1 with a sprint goal.
    \item Implemented the core learning flow: browse library $\rightarrow$ pick poem $\rightarrow$ learn chunk-by-chunk $\rightarrow$ test via text $\rightarrow$ score.
    \item Set up Git workflow: feature branches per issue, pull requests with descriptions, squash merges.
    \item Deployed MVP~v1 to the VPS.
    \item Conducted sprint review and retrospective.
\end{enumerate}

\textbf{Non-participation:} N/A (solo project --- all work done by Egor Zhukov).

% ============================================================
% SPRINT PLANNING
% ============================================================

\section{Sprint Planning}

\textbf{Repository:} \href{https://github.com/d13-l1t3/PoetryForYou}{github.com/d13-l1t3/PoetryForYou}

\textbf{Milestone:} \href{https://github.com/d13-l1t3/PoetryForYou/milestone/1}{Sprint 1 Milestone (link)}

\textbf{Sprint Goal:} \textit{By the end of this sprint, the customer can use the deployed bot (@PoetryforYou\_bot) to browse the poem library, pick a poem, learn it chunk-by-chunk via text input, and see a similarity score --- completing the core learning flow end-to-end.}

\textbf{Kanban Board:} \href{https://github.com/d13-l1t3/PoetryForYou/projects/1}{GitHub Projects Board (link)}

\subsection*{Story Points Tracking}

\begin{center}
\begin{tabular}{|l|r|}
\hline
\textbf{Metric} & \textbf{Value} \\
\hline
Total backlog (all issues) & 65 SP \\
\hline
Sprint 1 planned & 31 SP \\
\hline
Sprint 1 completed & 31 SP \\
\hline
Velocity & 31 SP/week \\
\hline
\end{tabular}
\end{center}

\subsection*{Sprint Backlog with Acceptance Criteria}

Issues pulled into Sprint~1 (Must-haves for the core flow):

\textbf{\#1 --- /start onboarding (2 SP)}\\
\texttt{GIVEN} the user opens the bot for the first time,\\
\texttt{WHEN} they send \texttt{/start},\\
\texttt{THEN} the bot greets them and asks to choose a language (ru/en/mix).

\textbf{\#2 --- Set difficulty level (2 SP)}\\
\texttt{GIVEN} the user has chosen a language,\\
\texttt{WHEN} the onboarding continues,\\
\texttt{THEN} the bot asks for difficulty (beginner/intermediate/advanced) and saves it.

\textbf{\#3 --- Browse library (5 SP)}\\
\texttt{GIVEN} the user is onboarded,\\
\texttt{WHEN} they send \texttt{/library},\\
\texttt{THEN} the bot shows a paginated list of poems filtered by language preference.

\textbf{\#4 --- Chunk-based learning (8 SP)}\\
\texttt{GIVEN} the user picks a poem from the library,\\
\texttt{WHEN} learning starts,\\
\texttt{THEN} the bot splits the poem into chunks (2--4 lines) and shows them one at a time with \texttt{/next} and \texttt{/stop} buttons.

\textbf{\#5 --- Text-based memory test (5 SP)}\\
\texttt{GIVEN} the user has memorized a chunk,\\
\texttt{WHEN} they press \texttt{/next},\\
\texttt{THEN} the bot hides the chunk and asks the user to type it from memory, returning a similarity score.

\textbf{\#15 --- Pause and resume (3 SP)}\\
\texttt{GIVEN} the user stops mid-poem with \texttt{/stop},\\
\texttt{WHEN} they return later and send \texttt{/review},\\
\texttt{THEN} the bot resumes from the last chunk they were learning.

\textbf{\#16 --- /help command (1 SP)}\\
\texttt{GIVEN} the user sends \texttt{/help},\\
\texttt{WHEN} the bot receives it,\\
\texttt{THEN} it displays all available commands with descriptions.

\textbf{\#7 --- Similarity score feedback (3 SP)}\\
\texttt{GIVEN} the user has typed or recited a chunk,\\
\texttt{WHEN} the bot compares their input to the original,\\
\texttt{THEN} it shows a percentage score; $\geq$80\% = success with points, $<$80\% = hint + retry.

\textbf{\#11 --- RU/EN interface (2 SP)}\\
\texttt{GIVEN} the user has set their language,\\
\texttt{WHEN} the bot sends any message,\\
\texttt{THEN} all UI text appears in the user's chosen language (RU or EN).

% ============================================================
% DEFINITION OF DONE
% ============================================================

\subsection*{Definition of Done}

An issue is ``Done'' when:
\begin{enumerate}
    \item Code is pushed to a feature branch and a Pull Request is created.
    \item PR description links the issue and describes the change.
    \item Code is reviewed (self-review for solo project: re-read the diff, check edge cases).
    \item PR is merged into \texttt{main} via squash merge.
    \item The feature works on the deployed bot (\texttt{@PoetryforYou\_bot}).
    \item Acceptance criteria pass when tested manually on the live bot.
\end{enumerate}

% ============================================================
% GIT WORKFLOW
% ============================================================

\section{Git Workflow}

\textbf{Branching:} One branch per issue, named \texttt{feature/\#N-short-description} (e.g., \texttt{feature/\#4-chunk-learning}).

\textbf{Commits:} Format: \texttt{[\#N] Short description} (e.g., \texttt{[\#4] Add chunk-based learning logic}).

\textbf{Pull Requests:} Each PR links the issue number, describes what was changed and how to test it. Merged via squash merge to keep \texttt{main} history clean.

\textbf{Review:} Self-review: re-read the full diff before merging, check for edge cases, test on the live bot. For solo project, every PR has at least one self-review comment.

\textbf{Deployment:} After merging to \texttt{main}, SSH into VPS, \texttt{git pull}, \texttt{docker compose up --build -d}.

% ============================================================
% SPRINT REVIEW
% ============================================================

\section{Sprint Review}

All 9 sprint backlog items (31 SP) were completed. The core learning flow works end-to-end:

\begin{enumerate}
    \item User sends \texttt{/start} $\rightarrow$ picks language $\rightarrow$ picks difficulty.
    \item User sends \texttt{/library} $\rightarrow$ browses poems $\rightarrow$ picks one.
    \item Bot shows chunk 1 $\rightarrow$ user memorizes $\rightarrow$ presses \texttt{/next} $\rightarrow$ types from memory $\rightarrow$ gets score.
    \item Repeats for all chunks $\rightarrow$ poem marked as completed $\rightarrow$ points earned.
    \item User can \texttt{/stop} mid-poem and resume later via \texttt{/review}.
    \item All UI text works in both RU and EN.
\end{enumerate}

% ============================================================
% RETROSPECTIVE
% ============================================================

\section{Sprint Retrospective}

\subsection*{What went well}
\begin{enumerate}
    \item The chunk-based learning flow works smoothly --- users don't get overwhelmed with long poems.
    \item Bilingual support (i18n) was implemented from the start, so no refactoring needed later.
    \item Docker deployment makes it easy to push updates: git pull + docker compose rebuild takes under 2 minutes.
\end{enumerate}

\subsection*{What didn't go well}
\begin{enumerate}
    \item Fuzzy matching for text comparison is too simplistic --- it doesn't handle line order or partial matches well.
    \item No automated tests yet --- all testing was done manually on the live bot.
    \item The Kanban board wasn't updated consistently during the sprint --- I tracked progress in my head instead.
\end{enumerate}

\subsection*{Action points}
\begin{enumerate}
    \item Improve the text comparison algorithm to handle partial matches and line reordering.
    \item Add unit tests for the core functions (chunk splitting, score calculation) in Sprint~2.
\end{enumerate}

% ============================================================
% MVP V1
% ============================================================

\section{MVP v1 Deployment}

\textbf{Live bot:} \href{https://t.me/PoetryforYou_bot}{@PoetryforYou\_bot}

\textbf{Repository:} \href{https://github.com/d13-l1t3/PoetryForYou}{github.com/d13-l1t3/PoetryForYou}

\textbf{What it does:} Onboarding (language + difficulty), poem library browsing, chunk-based learning, text-input memory testing with similarity score, pause/resume, bilingual UI (RU/EN), \texttt{/help} command.

% ============================================================
% LEARNINGS
% ============================================================

\section{What I Learned}

\begin{itemize}
    \item Chunk-based learning is the right approach --- it keeps the user engaged without overwhelming them.
    \item Text similarity checking needs more work: simple fuzzy matching doesn't handle line reordering or minor typos well.
    \item The Telegram reply-button UX makes the bot feel more conversational than raw text commands.
    \item Setting up Docker deployment early saved a lot of time --- now updating the live bot takes under 2 minutes.
\end{itemize}

% ============================================================
% A2 FEEDBACK ADDRESSED
% ============================================================

\section{A2 Feedback Addressed}

Based on the customer interview (Week~1) and A2 planning:
\begin{itemize}
    \item \textbf{Conversational UX} --- implemented reply buttons and natural dialogue flow, not a web-app-in-a-bot.
    \item \textbf{Core flow first} --- focused Sprint~1 entirely on browse $\rightarrow$ learn $\rightarrow$ test, no gamification yet.
    \item \textbf{Two languages} --- both RU and EN are supported from the start.
    \item \textbf{Chunk-based learning} --- customer said users shouldn't be overwhelmed; implemented 2--4 line chunks.
\end{itemize}

% ============================================================
% CURRENT PROGRESS + NEXT STEPS
% ============================================================

\section{Current Progress and Next Steps}

\subsection*{Where we are now}
MVP~v1 covers the core text-based learning flow. 9 out of 16 backlog items are done (31/65~SP). The bot is deployed and accessible.

\subsection*{Next steps (Sprint 2)}
\begin{enumerate}
    \item \textbf{Voice recitation} (\#6, 8~SP) --- integrate Whisper STT so users can recite via voice messages.
    \item \textbf{Online poem search} (\#8, 5~SP) --- let users search for poems beyond the built-in library.
    \item \textbf{Spaced repetition} (\#9, 5~SP) --- schedule reviews using SM-2 algorithm.
    \item \textbf{Progress dashboard} (\#10, 3~SP) --- show learning stats with \texttt{/progress}.
    \item Add unit tests for chunk splitting and score calculation.
\end{enumerate}

% ============================================================
% PRESENTATION
% ============================================================

\section{Presentation}

\href{https://PRESENTATION_SLIDES_LINK_HERE}{Presentation slides (link)} --- presented at lab on Week~3.

% ============================================================
% AI USAGE
% ============================================================

\section{AI Usage}

I used ChatGPT to help write acceptance criteria in GIVEN/WHEN/THEN format and to draft this report. The code, Git workflow, deployment, and all testing were done by me. AI was not used for writing application code this sprint.

\end{document}
